\documentclass[11pt,a4paper]{article}
\usepackage[T1]{fontenc}
\usepackage[utf8]{inputenc}
\usepackage{lmodern,amsmath,amssymb}
\usepackage[margin=25mm]{geometry}
\usepackage[hidelinks]{hyperref}
\hypersetup{pdftitle={Newton's cell action as a measure on cuts},pdfauthor={navstokgap research working notes}}
\newcommand{\tightlist}{\setlength{\itemsep}{0pt}\setlength{\parskip}{0pt}}
\setlength{\emergencystretch}{3em}
\setcounter{secnumdepth}{0}
\begin{document}
\hypertarget{newtons-cell-action-as-a-measure-on-cuts}{%
\section{Newton's cell action as a measure on
cuts}\label{newtons-cell-action-as-a-measure-on-cuts}}

\textbf{After refereeing (Fable, 2026-09-28).} Theorems 1--2, Corollary
3's arithmetic and the quotations were re-derived and accepted.
Corrections applied: Proposition 4 now names the Kullback--Leibler
divergence, the shift \(\delta x\), the phase \(K_\tau/\varepsilon\) and
the Cameron--Martin normalizer, and its link to Rivero 1998 is stated as
a counterpart; Corollary 5 separates a floor on a cell's action from a
floor on a cut's share (a factor \(\frac34\) for halving); §3 labels the
two-quarter-cut schedule a construction and adds the alternating reading
favoured by 𣃈必半; the Lévy--Ciesielski comparison names the right
quantity.

\textbf{Result, 2026-09-27 (atlas §1b, user remarks on arbitrary cuts
and on the rod).} Let a constant force \(F\) act on a body of mass \(m\)
over a cell of duration \(\tau\), and let \(\delta x\) be the chord
minus the parabola (the chord is the inertial comparison path with the
same endpoints). Its Galileo action

\[K_\tau=\frac m2\int_0^\tau\dot{\delta x}^2\,dt=\frac{F^2\tau^3}{24m}\]

(the \href{galileo-two-path-interference.md}{two-path note} writes the
same quantity as \((F/4v)A=\tau\Delta E/12\)) behaves as a measure on
the process of cutting the cell.

\begin{itemize}
\tightlist
\item
  \textbf{Theorem 1 (one cut).} A cut at fraction \(s\in(0,1)\) spends
  exactly \(3s(1-s)K_\tau\), and this share equals the Cameron--Martin
  energy \(\frac m2\int\dot\phi^2\) of the Schauder hat \(\phi\) that
  the cut inserts, of height the sagitta \(Fs(1-s)\tau^2/(2m)\).
\item
  \textbf{Theorem 2 (any cut sequence).} The hats of successive cuts are
  orthogonal in the energy \(\int\dot f\dot g\), so the spent shares
  add. Along any sequence of cuts, at any positions and in any order,
  the total spent equals \(K_\tau-\sum_iK_{\tau_i}\) over the current
  pieces, and it tends to \(K_\tau\) iff the mesh \(\max_i\tau_i\) tends
  to zero; the remainder is at most \(K_\tau\max_i\tau_i^2/\tau^2\).
\item
  \textbf{Corollary 3 (the rod and the Mohist schedules).} Halving only
  the remaining piece each day, the schedule of the stick in
  \emph{Zhuangzi} 33, spends \(\frac67K_\tau\), and so does one halving
  a day taken alternately from the front and the back. Removing a
  quarter from each end each day (two quarter-cuts a day, a construction
  motivated by the Mohist ``taking from front and back'') spends
  \(\frac{27}{28}K_\tau\). All three leave a definite share unspent
  forever, because they keep cut-off pieces of fixed length.
\item
  \textbf{Proposition 4 (weight of a cut).} Under the Euclidean path
  measure of a free particle with diffusion constant \(\hbar/m\) (Wiener
  measure), the force signal in one cut has Kullback--Leibler divergence
  exactly (spent share)\(/\hbar\). In real time the same share, divided
  by the action resolution, is the phase of the two-path comparison.
\item
  \textbf{Corollary 5 (a floor counts the cuts).} If every exhibited cut
  must spend at least an action \(\kappa>0\), at most \(K_\tau/\kappa\)
  cuts can be exhibited, whatever the schedule. For halving refinement,
  level \(n\) is exhibited iff \(\frac34K_\tau8^{-n}\ge\kappa\). The
  Planck paper's mark mesh
  \(\tau_*=(48z_{1-\epsilon}^2m\hbar/F^2)^{1/3}\) is the cell whose
  action is \(2z_{1-\epsilon}^2\hbar\); its own halving spends
  \(\frac32z_{1-\epsilon}^2\hbar\).
\end{itemize}

Theorems 1 and 2 are elementary, and their content lies in the reading
they give: Newton's refinement is a process that spends a fixed, finite
action, with the nested additivity of the Lévy--Ciesielski construction
of the Brownian bridge. Proposition 4 makes the correspondence exact:
the action is the Cameron--Martin energy of the force's shift,
decomposed over the hats. Corollary 5 is a consistency statement. It
says where a floor, if present, stops the cutting, and leaves the
necessity question of the
\href{principia-fifth-postulate.md}{fifth-postulate note} where it was.

\hypertarget{one-cut}{%
\subsection{1. One cut}\label{one-cut}}

Put \(C=F^2/(24m)\), so \(K_\tau=C\tau^3\). On \([0,\tau]\) the parabola
with acceleration \(F/m\) and the chord through its endpoints differ by

\[\delta x(t)=\frac F{2m}\,t(\tau-t),\qquad
\dot{\delta x}(t)=\frac Fm\Bigl(\frac\tau2-t\Bigr),\]

so
\(\frac m2\int_0^\tau\dot{\delta x}^2dt=\frac m2\cdot\frac{F^2}{m^2}\cdot\frac{\tau^3}{12}=C\tau^3\).

Cut at \(t_c=s\tau\). The polygon through the three points
\(0,s\tau,\tau\) differs from the chord by the hat \(\phi\) of height
\(h=\delta x(s\tau)=Fs(1-s)\tau^2/(2m)\), linear on \([0,s\tau]\) and on
\([s\tau,\tau]\). The polygon differs from the parabola, on each piece,
by the chord-minus- parabola of that piece:
\(\delta x=\phi+\delta x_1+\delta x_2\) with \(\delta x_j\) supported on
piece \(j\) and vanishing at its ends.

\emph{Orthogonality.} \(\dot\phi\) is constant on each piece and
\(\int_{\rm piece}\dot{\delta x_j}=0\), so
\(\int\dot\phi\,\dot{\delta x_j}=0\); the \(\delta x_j\) have disjoint
supports. Hence

\[K_\tau=\frac m2\int\dot\phi^2+K_{s\tau}+K_{(1-s)\tau}.\]

\emph{The hat's energy.}
\(\frac m2\int\dot\phi^2=\frac m2h^2\bigl(\frac1{s\tau}+\frac1{(1-s)\tau}\bigr) =\frac{mh^2}{2s(1-s)\tau}=\frac{F^2s(1-s)\tau^3}{8m}=3s(1-s)C\tau^3.\)
The same number follows from
\(\tau^3=(s\tau)^3+((1-s)\tau)^3+3s(1-s)\tau^3\), which is the algebraic
form of Theorem 1. The share is largest for halving, \(\frac34K_\tau\).
\(\square\)

\hypertarget{any-cut-sequence}{%
\subsection{2. Any cut sequence}\label{any-cut-sequence}}

Each later cut acts inside one current piece, and Theorem 1 applies to
that piece with its own \(s\). The new hat is supported in the piece and
vanishes at its ends; every earlier hat is linear there. The argument of
§1 gives orthogonality to all earlier hats and to the other pieces'
remainders. After any finite set of cuts, with pieces \(\tau_i\)
(\(\sum\tau_i=\tau\)),

\[K_\tau=\sum_{\rm cuts}\frac m2\int\dot\phi_{\rm cut}^2+\sum_iK_{\tau_i},\qquad
\sum_iK_{\tau_i}=C\sum_i\tau_i^3\le C\tau\max_i\tau_i^2 .\]

The remainder vanishes iff the mesh does
(\(\sum_i\tau_i^3\ge(\max_i\tau_i)^3\)). The positions and the order
enter only through which pieces are cut, never through the total.
\(\square\)

This is the Schauder (Lévy--Ciesielski) expansion of \(\delta x\) in the
Cameron--Martin space \(H^1_0[0,\tau]\), adapted to an arbitrary nested
cut sequence; for dyadic cuts it is the textbook expansion
(\href{https://doi.org/10.1090/S0002-9947-1961-0132591-2}{Ciesielski
1961}; metadata). The \href{reachable-cut-composition.md}{reachability
note} has the same cubic law for the area \(2F^2t^3/(3m)\) of the
reachable lens of a cut. The bridge's own conditional variance shows the
same nested additivity with a quadratic law: for diffusion constant
\(D\), \(\int_0^L{\rm Var}(X_u\mid{\rm ends})\,du=DL^2/6\), and a cut at
fraction \(s\) removes \(Ds(1-s)L^2/3\) of it.

\hypertarget{the-rod-and-the-two-mohist-schedules}{%
\subsection{3. The rod and the two Mohist
schedules}\label{the-rod-and-the-two-mohist-schedules}}

\emph{Zhuangzi} 33 (dialecticians' list,
\href{../docs/classics/Zhuangzi_33_Tianxia_zh_wikisource.md}{text}):
一尺之捶，日取其半，萬世不竭, ``a one-foot stick, each day take half, in
ten thousand generations it is not exhausted.'' Mohist Canon B (B60 in
Graham's numbering;
\href{../docs/classics/Mozi_Canons_B_JingXia_zh_wikisource.md}{canon}):
非半弗𣃈，則不動，說在端, ``without halving there is no cutting and no
moving; the explanation lies in the point (端)'', with the explanation
(\href{../docs/classics/Mozi_Canons_B_Explanations_JingShuoXia_zh_wikisource.md}{text})
前則中無爲半，猶端也。前後取，則端中也 (working gloss: ``going forward,
at the middle nothing serves as a half: it is like a point; taking from
front and back, the point is in the middle''), which continues 𣃈必半,
``cutting must be by halves''.

\emph{One-ended schedule.} Day \(k\) halves the remaining piece, of
length \(\tau2^{-(k-1)}\), and spends \(\frac34C\tau^38^{-(k-1)}\). The
total is \(\frac34\cdot\frac87K_\tau=\frac67K_\tau\); the pieces taken
away keep \(\sum_{k\ge1}8^{-k}K_\tau=\frac17K_\tau\).

\emph{Alternating schedule} (the reading favoured by 𣃈必半). One
halving a day, taking the half alternately from the front and from the
back. Each day is one halving of the remaining piece, so the spending is
that of the one-ended schedule, \(\frac67K_\tau\); the remainder
converges to the point \(\frac\tau2+\frac\tau8+\dots=\frac{2\tau}3\).

\emph{Two quarter-cuts a day} (a construction motivated by 前後取; the
text does not fix the fractions). Day \(k\) cuts the remaining piece
\(r=\tau2^{-(k-1)}\) at \(s=\frac14\) and \(s=\frac34\), keeping the
middle half. The two cuts spend
\(C r^3\bigl(1-2\cdot\frac1{64}-\frac18\bigr)=\frac{27}{32}Cr^3\)
(equivalently \(\frac9{16}+\frac9{32}\) by two applications of Theorem
1), and the total is
\(\frac{27}{32}\cdot\frac87K_\tau=\frac{27}{28}K_\tau\), with
\(\frac1{28}K_\tau\) kept in the pieces taken away. The contrast
\(\frac67\) against \(\frac{27}{28}\) is between one half-cut and two
quarter-cuts a day, whichever ends are cut.

All three schedules converge to a point, the Mohist 端: the end, the
point \(2\tau/3\), and the midpoint. If 中 in 端中 means the midpoint,
the explanation fits the symmetric schedule, while 𣃈必半 fits the
alternating one. What each author is committed to, and what the theorem
does with it:

\begin{itemize}
\tightlist
\item
  The dialecticians hold that the halving never ends. Theorem 2 agrees
  for the figure and adds a price: their schedule spends only
  \(\frac67\) of the action, since it never cuts again what it took
  away.
\item
  The Mohists hold that the halving ends at a point. Theorem 2 agrees
  that each schedule converges to a point, and the point carries no
  action. The unhalvable 端 as a least magnitude is a reading the
  theorem leaves aside, as the \href{planck-gap-paper.md}{Planck paper}
  already records.
\item
  Newton's scholium closing Book I, Section I, takes the geometers' side
  (divisibility without end). Theorem 2 is that side made quantitative:
  refinement exhausts the action only when every piece is cut again.
\end{itemize}

The parallel with Zeno's dichotomy (Aristotle, \emph{Physics} VI.9) is
recorded as convergence; transmission would need evidence this note does
not have.

\hypertarget{the-weight-of-a-cut}{%
\subsection{4. The weight of a cut}\label{the-weight-of-a-cut}}

Take the Euclidean path measure of a free particle with diffusion
constant \(\hbar/m\): Brownian motion \(X\) with \(E[X_t^2]=\hbar t/m\),
pinned at the cell's endpoints. Given the endpoints of a piece of length
\(L\), the value at fraction \(s\) is Gaussian with variance
\(\hbar s(1-s)L/m\), independent of the finer hats and of the other
pieces, by the Markov property of the bridge (the Lévy construction, for
any nested sequence). A force \(F\) shifts the mean of that value by the
sagitta \(h=Fs(1-s)L^2/(2m)\). The Kullback--Leibler divergence of the
shifted from the free law (the expected log-likelihood ratio, which here
also equals the ratio evaluated at the shift) is

\[\frac{h^2}{2\,\hbar s(1-s)L/m}=\frac{mh^2}{2s(1-s)L\,\hbar}
=\frac{3s(1-s)CL^3}{\hbar},\]

the spent share of Theorem 1 divided by \(\hbar\). By the chain rule the
divergences add over cuts, and along a sequence with mesh tending to
zero they sum to \(K_\tau/\hbar\), the Cameron--Martin exponent for the
shift \(\delta x\). The Cameron--Martin density factorizes over cuts,
each contributing the normalizing factor \(e^{-3s(1-s)CL^3/\hbar}\).
This is the Euclidean counterpart of the real-time two-path phase
\(K_\tau/\varepsilon\) of the two-path note (the same number with
\(\varepsilon\) in the role of \(\hbar\)), and of the product-over-cuts
structure that §7 of the
\href{principia-fifth-postulate.md}{fifth-postulate note} reads in the
time-step partition of
\href{https://arxiv.org/abs/quant-ph/9803035}{Rivero 1998}. For
\(\hbar\to0\) every cut, however small, is decisive; for \(\hbar>0\) the
cuts whose share falls below \(\hbar\) carry little evidence (under one
nat). \(\square\)

\hypertarget{a-floor-counts-the-cuts}{%
\subsection{5. A floor counts the cuts}\label{a-floor-counts-the-cuts}}

If every cut that a record exhibits must spend at least \(\kappa>0\),
Theorem 2 gives at most \(K_\tau/\kappa\) such cuts, for every schedule.
For halving refinement the \(2^n\) cuts of level \(n\) each spend
\(\frac34K_\tau8^{-n}\), so the exhibited levels are those with
\(8^n\le 3K_\tau/(4\kappa)\).

Two floors must be kept apart: \(\kappa\) here floors the share a cut
spends, \(3s(1-s)K_L\), while the thresholds in the other notes floor
the action of a cell, \(K_\tau\); for halving the two differ by the
factor \(\frac34\). The Planck paper's eq. (2) is
\(K_\tau\ge2z_{1-\epsilon}^2\hbar\): \(\tau_*\) is the cell with
\(K_{\tau_*}=2z_{1-\epsilon}^2\hbar\), whose halving spends
\(\frac32z_{1-\epsilon}^2\hbar\). The round-3 resolution threshold
\(\tau\ge(24m\varepsilon d_{n,\eta}/F^2)^{1/3}\) is
\(K_\tau\ge\varepsilon d_{n,\eta}\), that is
\(\kappa=\frac34\varepsilon d_{n,\eta}\) in cut-share form. Corollary 5
thus places a floor, when one is given, at a definite depth of the cut
process; it supplies no floor.

\hypertarget{consequence-for-state}{%
\subsection{6. Consequence for STATE}\label{consequence-for-state}}

Atlas §1b: the Newton row of ``cuts at any position'' is proved here
(Theorems 1--2, Corollary 3), with the Lévy--Ciesielski identification
and the Euclidean weight of a cut (Proposition 4). The necessity
question is unchanged. On the gauge side of §1b the free-field defect of
a cut at fraction \(s\) is (5) scaled by \(4s(1-s)\) in its size bound
(\href{series-parallel-gauge-refinement.md}{Corollary 2\(_s\)}); the
\(U(1)\) Theorem 5 holds for every cut fraction (Corollary 5\(_s\)
there).
\end{document}
